\documentclass[11pt]{article}
\usepackage{setspace}
\usepackage{geometry}
\geometry{margin=1in}
\setstretch{1.15}

\begin{document}

% --- Applicant Information ---
\noindent
\textbf{Zhongli Wang} \\
Phone: +49 160 9541 6015 \\
Email: nkwangzhongli@gmail.com \\
Website: https://sites.google.com/view/zhongli-wang \\
Bielefeld, Germany \\

\vspace{0.5cm}

% --- Recipient Information ---
\noindent
\textbf{Prof.\ Dr.\ Christine Laudenbach} \\
Leibniz Institute for Financial Research SAFE \\
House of Finance, Goethe University Frankfurt \\
Frankfurt am Main, Germany \\

\vspace{0.5cm}

% --- Subject Line ---
\noindent
\textbf{Application for Postdoctoral Researcher in Financial Intermediation / Household Finance} \\

\vspace{0.3cm}

\noindent Dear Prof.\ Laudenbach and Members of the Selection Committee,\\[0.2cm]

I am writing to apply for the Postdoctoral Researcher position in \textit{Financial Intermediation or Household Finance} at SAFE. 
I am currently completing my Ph.D.\ in Economics at Bielefeld University and Università Cattolica del Sacro Cuore within the Marie Skłodowska--Curie EPOC programme. 
My research combines applied microeconomics, quantitative modeling, and computational methods to study how incentives, information frictions, and institutional design shape market outcomes. 
While my dissertation is rooted in industrial organization and digital markets, my methodological focus and several projects directly connect to core questions in finance---in particular, how institutions and households respond to constraints, risks, and policy-relevant changes in the information environment. 
SAFE's mission to strengthen a resilient and sustainable financial system, together with its strong research community at the House of Finance, makes it an ideal environment to develop this agenda.

A central theme of my work is the analysis of strategic behavior under frictions and heterogeneity. 
In my job market paper, I develop and estimate a dynamic agent-based model in which competing firms adjust capacity investment, pricing, and service quality under network effects and congestion. 
Methodologically, this work translates naturally to financial intermediation settings where balance-sheet constraints, funding conditions, and interconnectedness generate non-linear dynamics and amplification mechanisms. 
Building on this foundation, I plan to develop research on how intermediary incentives and market design shape the allocation of risk and liquidity, and how these mechanisms interact with regulation and the adoption of new financial technologies.

In parallel, I have pursued research on financial market behavior, transparency, and data-driven risk detection. 
In \textit{Polytope Fraud Theory} (\textit{International Review of Financial Analysis}, 2025), I apply supervised machine learning to identify structured irregularities in corporate financial statements, focusing on firms targeted by activist short sellers in U.S.\ markets. 
The project links monitoring incentives and institutional features to observable anomalies and provides a unified framework for investor protection. 
This experience strengthened my ability to work with granular firm-level data, build reproducible empirical pipelines, and translate quantitative evidence into economically meaningful mechanisms---skills relevant to both household finance (e.g., financial literacy, misperceptions, disclosure) and intermediation (e.g., screening, risk-taking, discipline, and market integrity).

I am also inspired by the household finance perspective in \textit{Rethinking the Stock Market Participation Puzzle: A Qualitative Approach} (with Kamila Duraj, Daniela Grunow, Michael Haliassos, and Stephan Siegel). 
The paper's emphasis on beliefs, constraints, and lived decision processes motivates a natural extension in my own work: combining qualitative or survey-based evidence on participation barriers with a disciplined theoretical framework that allows for heterogeneous agents, belief updating, and non-standard constraints. 
My experience with agent-based and simulation-based models positions me to contribute to such a research direction by translating empirically grounded mechanisms into tractable heterogeneous-agent structures that generate testable implications for participation, portfolio choice, and policy interventions.

I also coauthored \textit{Agentic Workflows for Economic Research: Design and Implementation} (arXiv, 2025), which develops a framework for integrating large-language-model tools into empirical workflows with explicit attention to documentation, reproducibility, and governance. 
This line of work is relevant for modern empirical finance, where researchers increasingly combine administrative or commercial datasets with unstructured data sources and complex pipelines. 
At SAFE, I would be keen to contribute to transparent and scalable workflows for handling firm reports, disclosures, regulatory texts, or other unstructured inputs that can inform research on intermediation, household expectations, and policy design.

SAFE offers a uniquely suitable environment for my next research stage. 
I am particularly motivated by the Institute's combination of frontier academic research, policy relevance, and interdisciplinary exchange across economics, law, and political science. 
I would welcome the opportunity to engage with SAFE's seminar culture, contribute to collaborative projects, and participate in the Institute's research and transfer activities. 
If opportunities arise, I would also be happy to contribute to teaching and student supervision at Goethe University Frankfurt, especially in applied econometrics, computational methods, or empirical finance.

Thank you so much for considering my application! 
I would be delighted to discuss how my research experience and quantitative skill set can contribute to SAFE's work on financial intermediation and household finance.

\vspace{6mm}

\noindent Sincerely,\\[15mm]
\noindent Zhongli Wang

\end{document}